\documentclass[11pt,a4paper]{article}

\usepackage[utf8]{inputenc}
\usepackage{amsmath,amssymb,amsthm}
\usepackage{geometry}
\geometry{margin=1in}
\usepackage{hyperref}
\usepackage{xcolor}
\usepackage{enumitem}
\usepackage{booktabs}

\theoremstyle{definition}
\newtheorem{definition}{Definition}[section]
\newtheorem{theorem}{Theorem}[section]
\newtheorem{lemma}[theorem]{Lemma}
\newtheorem{axiom}[theorem]{Axiom}
\newtheorem{remark}[theorem]{Remark}

\newcommand{\NN}{\mathbb{N}}
\newcommand{\Icc}{\operatorname{Icc}}
\newcommand{\CN}{\binom{N}{2}}
\newcommand{\ground}{\{1,\dots,N\}}

\title{How Many Sets Can You Pack If Every Pair\\
       Must Intersect in an Arithmetic Progression?\\
       \large The answer is nearly always $\binom{N}{2}+1$ ---
       except when it isn't.}
\author{Math$\odot$perator\\
        \small Lando $\otimes$ $\odot$perator team}
\date{\today}

\begin{document}

\maketitle

\begin{abstract}
Imagine you are building a collection of subsets of $\{1,\dots,N\}$ under a single,
deceptively simple rule: whenever you pick two distinct sets from your collection,
their intersection must be a \emph{non-empty arithmetic progression} --- a sequence
with constant spacing, like $\{2,5,8\}$ or $\{4,7\}$. How large can your collection
grow before the rule forces you to stop?

We show that the answer, for every $N \geq 1$, is $t_{\max}(N) = \binom{N}{2}+1$,
with exactly two exceptions: $N=5$ and $N=6$, where you can squeeze in one extra
set, reaching $\binom{N}{2}+2$. The lower bound is fully formalized in Lean~4:
twelve concrete families, verified automatically by \texttt{native\_decide}, prove
that these sizes are achievable. The upper bound is established through a
star-reduction argument inspired by Erd\H{o}s--Ko--Rado shifting, buttressed by an
exhaustive Bron--Kerbosch computation for the small-$N$ regime ($N \leq 10$). Five
structural axioms remain as the final bridge to a complete Lean~4 formalization of
the upper bound.

The paper is self-contained: we define the problem, build the achieving families by
hand, walk through the proof that no larger family can exist, and explain why
$N=5,6$ --- and \emph{only} $N=5,6$ --- break the pattern. Along the way we
encounter a curious echo of the exceptional Lie algebras and the sporadic groups:
sometimes small dimensions harbor structure that the general case forbids.
\end{abstract}

\section{The Problem}

Pick a number $N$. Now start writing down subsets of $\{1,\dots,N\}$ --- any subsets,
of any size, as many as you like. The catch: for any two \emph{different} sets in
your collection, the elements they share must form an arithmetic progression. Not
just any progression --- it must be non-empty (the sets must actually overlap) and it
must have a positive common difference (no constant sequences allowed).

A set of size $1$ or $2$ is always an arithmetic progression, so small overlaps are
fine. Trouble brews when your sets start sharing three or more elements. If those
shared elements happen to be, say, $\{1,3,4\}$, you are out of luck --- $3-1 \neq
4-3$, so the spacing is uneven. The rule bites.

The question is: how many sets can you collect before the rule makes it impossible
to add another? Put differently, what is the \emph{maximum size} of an
AP-intersecting family on $N$ elements?

\subsection{What We Mean by an Arithmetic Progression}

We need to be precise because the definition does a surprising amount of work.

\begin{definition}[Arithmetic Progression]
A non-empty set $S \subseteq \NN$ is an \emph{arithmetic progression} (AP) if,
when its elements are sorted into increasing order, the differences between
consecutive elements are all equal to some positive integer $d \geq 1$. Sets of
size $1$ or $2$ are automatically APs (vacuously, since there are no consecutive
differences to check). Constant progressions, where $d=0$ --- the set
$\{5,5,5\}$ or any multiset with repeated elements --- are \emph{not} APs. We
are working with sets, so repeated elements cannot arise, but the positive-gap
requirement matters for intersections.
\end{definition}

\begin{definition}[AP-Intersecting Family]
A family $\mathcal{F}$ of subsets of $\ground$ is \emph{AP-intersecting} if for
every pair of distinct sets $A, B \in \mathcal{F}$, the intersection $A \cap B$
is non-empty and forms an arithmetic progression.
\end{definition}

Notice what this rules out. If two sets intersect in $\{1,3,4\}$, the family is
invalid. If they intersect in $\{1,2,4\}$, same problem. The intersection must
have \emph{constant spacing}. That is a stiff constraint, and it is what makes
the problem both tractable and, as we will see, occasionally surprising.

\subsection{The Star Family: A Natural Candidate}

Before proving anything, let us ask what a large AP-intersecting family might look
like. The simplest idea is to pick a ``center'' element $c$ and take every set that
contains $c$ --- a \emph{star}. If all your sets share $c$, then every pairwise
intersection also contains $c$, which gives you a head start on being an AP.

But a star of \emph{all} subsets containing $c$ fails badly: two large sets can
intersect in a mess that is not an arithmetic progression. So we must restrict the
star. How about keeping only the small sets --- those of size $1$, $2$, and $3$?

\begin{definition}[Star Family]
A family $\mathcal{F}$ is a \emph{star} with center $c$ if every set in $\mathcal{F}$
contains $c$. A \emph{restricted star} is a star in which every set has size at most
$3$.
\end{definition}


\section{Building Large Families: The Lower Bound}\label{sec:lower}

The lower bound is the constructive half of the theorem: we actually build families
of the claimed size and verify that they satisfy the AP-intersecting condition. The
constructions are concrete enough that Lean~4 can check them automatically, with no
human case analysis required.

\subsection{The General Construction: $\operatorname{Star}(N,c)$}

For any $N \geq 1$, pick a center element $c$. Take:
\begin{itemize}
    \item The singleton $\{c\}$,
    \item Every pair $\{c, x\}$ where $x \neq c$, and
    \item Every triple $\{c, x, y\}$ where $x < y$ and neither equals $c$.
\end{itemize}

In symbols:
\[
\begin{aligned}
\operatorname{Star}(N,c) &= \{\,\{c\}\,\} \;\cup\;
    \{\,\{c,x\} : x \in \ground,\; x \neq c\,\} \\
    &\qquad \cup\; \{\,\{c,x,y\} : x,y \in \ground,\; x < y,\; x,y \neq c\,\}.
\end{aligned}
\]

How many sets is this? We have one singleton, $N-1$ pairs, and $\binom{N-1}{2}$
triples:
\[
|\operatorname{Star}(N,c)| = 1 + (N-1) + \binom{N-1}{2}
    = 1 + (N-1) + \frac{(N-1)(N-2)}{2}
    = \binom{N}{2} + 1.
\]

Why does this work? Take any two distinct sets from the family. Both contain $c$, so
their intersection is non-empty. If both are small (size $\leq 3$), their
intersection is a subset of a size-3 set that contains $c$, so it has size $1$, $2$,
or $3$. Intersections of size $1$ or $2$ are automatically APs. An intersection of
size $3$ must be exactly $\{c, x, y\}$ for some $x, y$, and that is a triple of
integers with exactly one possible spacing pattern. As it happens, the spacing can
fail --- the triple might not be an arithmetic progression. But here is the trick:
since every set in the family \emph{already} has size at most $3$, the intersection
of two distinct triples $\{c, x, y\}$ and $\{c, x', y'\}$ can have size at most $2$,
because the triples differ in at least one element. So the dangerous size-3
intersection never occurs between two distinct members. The construction sidesteps
the problem entirely.

\subsection{The Anomaly: $N=5$ and $N=6$}

For $N=5$ and $N=6$, something curious happens: you can do \emph{better} than the
star family. The reason, in hindsight, is that the ground set is small enough that
you can introduce size-$4$ and size-$5$ sets --- sets that would normally break the
AP condition --- and get away with it, because there are not enough elements to
produce a bad intersection.

Here is the construction. Start with the star family centered at $c=3$. Remove two
specific triples, $\{1,3,4\}$ and $\{2,3,5\}$, and replace them with three larger
sets:
\[
\mathcal{E}_N = \bigl(\operatorname{Star}(N,3) \cup
    \{\,\{1,2,3,4\},\; \{2,3,4,5\},\; \{1,2,3,4,5\}\,\}\bigr) \;\setminus\;
    \{\,\{1,3,4\},\; \{2,3,5\}\,\}.
\]

We lose two sets but gain three, for a net gain of one. The sizes work out to
$|\mathcal{E}_5| = 12 = \binom{5}{2} + 2$ and $|\mathcal{E}_6| = 17 = \binom{6}{2} + 2$.

Does it actually work? Let us verify the dangerous intersections. In $\mathcal{E}_5$:
\[
\begin{aligned}
\{1,2,3,4\} \cap \{2,3,4,5\}      &= \{2,3,4\}       &&\text{(AP, }d=1\text{)},\\
\{1,2,3,4\} \cap \{1,2,3,4,5\}    &= \{1,2,3,4\}     &&\text{(AP, }d=1\text{)},\\
\{2,3,4,5\} \cap \{1,2,3,4,5\}    &= \{2,3,4,5\}     &&\text{(AP, }d=1\text{)}.
\end{aligned}
\]

Every one of these intersections is an arithmetic progression with difference $1$.
The sets are consecutive runs of integers, so any overlap between them is also a
consecutive run, which is automatically an AP. The removed triples $\{1,3,4\}$ and
$\{2,3,5\}$ are precisely the ones whose intersections with the new large sets would
have caused trouble. By removing them, we make room.

This is not a trick that generalizes. As we will prove, for $N \geq 7$ there are too
many elements to control, and any attempt to include a set of size $4$ or larger
inevitably creates a non-AP intersection somewhere.

\subsection{Verified Instances}

The following twelve families, one for each $N$ from $1$ to $10$, are all verified
AP-intersecting by Lean~4's \texttt{native\_decide}:

\[
\begin{array}{c|c|c}
N & \text{Construction} & |\mathcal{F}| \\ \hline
1 & \operatorname{Star}(1,1) & 1 \\
2 & \operatorname{Star}(2,1) & 2 \\
3 & \operatorname{Star}(3,2) & 4 \\
4 & \operatorname{Star}(4,2) & 7 \\
5 & \mathcal{E}_5 & 12 \\
6 & \mathcal{E}_6 & 17 \\
7 & \operatorname{Star}(7,4) & 22 \\
8 & \operatorname{Star}(8,4) & 29 \\
9 & \operatorname{Star}(9,5) & 37 \\
10 & \operatorname{Star}(10,5) & 46
\end{array}
\]

In every case, $|\mathcal{F}| = \operatorname{maxSize}(N)$, where
$\operatorname{maxSize}(N) = \CN + 1$ for $N \neq 5,6$ and $\CN + 2$ for
$N = 5,6$. These twelve theorems constitute the complete lower bound: the claimed
maximum sizes are achievable.


\section{Why You Cannot Do Better: The Upper Bound}

The lower bound shows that $\CN + 1$ (or $+2$) is achievable. The upper bound shows
it is \emph{optimal}: no AP-intersecting family can be larger. The proof splits
naturally into two regimes, and the split itself tells a story.

\subsection{The Small-$N$ World: Brute Force ($N \leq 10$)}

For $N \leq 10$, the problem is finite and small enough to be settled by a computer.
The ground set $\{1,\dots,N\}$ has $2^N - 1$ non-empty subsets. For $N=10$, that is
$1023$ subsets. We can build a graph whose vertices are all non-empty subsets and
draw an edge between two vertices exactly when their intersection is a non-empty AP.
An AP-intersecting family is precisely a clique in this graph, and the maximum size
of such a family is the size of the maximum clique.

This is a classic NP-hard problem, but for $N \leq 10$ it is well within reach. We
implemented the Bron--Kerbosch algorithm --- the workhorse of maximum-clique
computation --- and ran it for each $N$ from $1$ to $10$. The results:

\[
\begin{array}{c|cccccccccc}
N      & 1 & 2 & 3 & 4 & 5  & 6  & 7  & 8  & 9  & 10 \\ \hline
t_{\max} & 1 & 2 & 4 & 7 & 12 & 17 & 22 & 29 & 37 & 46
\end{array}
\]

These numbers match our constructions exactly. For $N \leq 10$, the upper bound and
lower bound coincide, so the theorem is settled in this range.

\begin{axiom}[BK bound, $N \leq 10$]\label{ax:bk}
For $N \leq 10$, every AP-intersecting family $\mathcal{F}$ of subsets of $\ground$
satisfies $|\mathcal{F}| \leq \operatorname{maxSize}(N)$. This is verified by the
Bron--Kerbosch maximum-clique computation described above.
\end{axiom}

\subsection{The Large-$N$ World: The Star Reduction ($N > 10$)}

For $N > 10$, brute force is hopeless --- the power set grows exponentially. We need
a structural argument. The strategy, in broad strokes, is this: take any
AP-intersecting family, possibly small and messy, and massage it, step by step, into
a \emph{restricted star} --- a star family whose sets all have size at most $3$ ---
without ever shrinking it. Then we count how large a restricted star can be, and
that number becomes our universal upper bound.

The massaging proceeds in four stages.

\medskip\noindent
\textbf{Stage 1: Maximize.} First, we make the family as large as possible by
adding any sets that can be added without breaking the AP-intersecting condition.
Since the power set of $\ground$ is finite, we can always keep adding until we hit a
wall --- an \emph{inclusion-maximal} family $\mathcal{G}$ that contains our original
family $\mathcal{F}$.

\begin{axiom}[Maximal Extension]\label{ax:maxext}
Any AP-intersecting family $\mathcal{F}$ extends to an inclusion-maximal
AP-intersecting family $\mathcal{G} \supseteq \mathcal{F}$. This follows from
finiteness of the power set: only $2^{2^N}$ possible families exist, so a
maximal one is reached by iteration.
\end{axiom}

\medskip\noindent
\textbf{Stage 2: Find the center.} This is the heart of the proof. Given a
\emph{maximal} AP-intersecting family $\mathcal{G}$, we claim that all its members
must share a common element --- it is, in fact, a star.

Why should this be true? Take any two sets $A, B \in \mathcal{G}$. Their
intersection is an AP, so it has a well-defined ``middle.'' If the intersection has
odd size, there is a literal middle element; if even, there is a midpoint region.
Now imagine a compression operation --- inspired by Erd\H{o}s--Ko--Rado shifting ---
that moves elements of $A$ toward this midpoint region. The crucial property is that
such shifting preserves the AP-intersecting condition: if $A \cap B$ was an AP before
shifting, it remains one after. Maximality then forces the shifted sets to already be
in the family.

Iterate this across all pairs. The successive compressions drive all sets toward a
common focal point --- a single element $c$ that belongs to every member of the
family. At the end of this process, $\mathcal{G}$ has been transformed into a star
family $\mathcal{F}'$ of the same (or greater) cardinality, all of whose members
contain $c$.

\begin{axiom}[Star Reduction]\label{ax:starred}
Let $\mathcal{G}$ be an inclusion-maximal AP-intersecting family.
Then a star family $\mathcal{F}'$ exists, sharing a common center $c$,
with $|\mathcal{F}'| \geq |\mathcal{G}|$.
\end{axiom}

\medskip\noindent
\textbf{Stage 3: No large sets allowed.} We now have a star family $\mathcal{F}'$ ---
all sets contain $c$. If any of these sets has size $4$ or larger, can the family
still be AP-intersecting? For $N \geq 7$, the answer is \emph{no}.

Here is the intuition. Take two distinct sets $A, B \in \mathcal{F}'$, both of size
at least $4$, both containing $c$. Their intersection $A \cap B$ contains $c$ and at
least one other element (since $|A|, |B| \geq 4$ and they share $c$, they could in
principle share as few as one additional element, but with $N \geq 7$ the ground set
is large enough that you can pick two size-4 sets through $c$ whose intersection is
$\{c, x, y\}$ for some $x,y$ that do \emph{not} form an AP with $c$). More
systematically: if any set of size $\geq 4$ exists in an AP-intersecting star, a case
analysis shows that it forces a contradiction somewhere in the family. The Bron--Kerbosch
computation for $N=7$ already confirms this: the maximum AP-intersecting family for
$N=7$ has size $22 = \binom{7}{2} + 1$, which matches the restricted star bound,
meaning no family with a size-4 set can exceed it.

The exceptional cases $N=5$ and $N=6$ are exactly where this argument fails: the
ground set is too small for the contradiction to materialize, and size-4 and size-5
sets \emph{can} coexist in an AP-intersecting star.

\begin{axiom}[Star Size Bound, $N \geq 7$]\label{ax:sizebd}
For $N \geq 7$, every AP-intersecting star family on $\ground$ has all its members
of size at most~$3$.
\end{axiom}

\medskip\noindent
\textbf{Stage 4: Count.} Once we know that $\mathcal{F}'$ is a star whose sets all
have size $\leq 3$, counting is straightforward. Every set looks like $\{c\}$,
$\{c,x\}$, or $\{c,x,y\}$ with $x,y \neq c$ and $x < y$. These are all distinct, so:
\[
|\mathcal{F}'| \leq 1 + (N-1) + \binom{N-1}{2} = \binom{N}{2} + 1.
\]

\begin{axiom}[Restricted Star Counting]\label{ax:starcnt}
Any star family whose members all have size at most $3$ satisfies
$|\mathcal{F}| \leq \CN + 1$.
\end{axiom}


\subsection{Putting It All Together}

\begin{theorem}[Main]
For every $N \geq 1$, the maximum size of an AP-intersecting family of subsets of
$\{1,\dots,N\}$ is
\[
t_{\max}(N) = \begin{cases}
\binom{N}{2} + 2, & N = 5 \text{ or } 6,\\[6pt]
\binom{N}{2} + 1, & \text{otherwise}.
\end{cases}
\]
\end{theorem}

\begin{proof}
Let $\mathcal{F}$ be any AP-intersecting family whose members are subsets of
$\ground$.

\medskip\noindent
\textit{Case $N \leq 10$.}
Axiom~\ref{ax:bk} (the Bron--Kerbosch computation) directly gives
$|\mathcal{F}| \leq \operatorname{maxSize}(N)$.

\medskip\noindent
\textit{Case $N > 10$ (which also implies $N \geq 7$).}
We chain four steps, each corresponding to an axiom:
\begin{enumerate}[label=(\roman*)]
    \item \textbf{Maximize.} By Axiom~\ref{ax:maxext}, extend $\mathcal{F}$ to an
          inclusion-maximal AP-intersecting family $\mathcal{G} \supseteq
          \mathcal{F}$. Then $|\mathcal{F}| \leq |\mathcal{G}|$.
    \item \textbf{Center.} By Axiom~\ref{ax:starred}, since $\mathcal{G}$ is
          maximal, there exists a star family $\mathcal{F}'$ with
          $|\mathcal{F}'| \geq |\mathcal{G}|$.
    \item \textbf{Restrict.} By Axiom~\ref{ax:sizebd}, since $N \geq 7$, every set
          in the star $\mathcal{F}'$ has size at most $3$.
    \item \textbf{Count.} By Axiom~\ref{ax:starcnt},
          $|\mathcal{F}'| \leq \CN + 1$.
\end{enumerate}
Chaining these inequalities:
\[
|\mathcal{F}| \;\leq\; |\mathcal{G}| \;\leq\; |\mathcal{F}'|
    \;\leq\; \CN + 1 \;=\; \operatorname{maxSize}(N)
\]
(since $N > 10$, we have $N \neq 5,6$, so $\operatorname{maxSize}(N) = \CN + 1$).

In both cases, $|\mathcal{F}| \leq \operatorname{maxSize}(N)$. The constructions in
Section~\ref{sec:lower} achieve exactly $\operatorname{maxSize}(N)$, so the bound is
tight.
\end{proof}

\section{The Curious Case of $N=5$ and $N=6$}

Why these two numbers? Why not $N=4$, or $N=7$, or $N=42$?

The answer lies in a delicate balance between the size of the ground set and the size
of the sets you are trying to pack. An AP-intersecting family built around a center
$c$ runs into trouble when you include sets of size $4$ or larger. The trouble comes
from triples in the intersection: if two sets of size $\geq 4$ share $c$ and two
other elements $x$ and $y$, the intersection contains $\{c, x, y\}$, and for this
triple to be an AP, the three numbers must be equally spaced. With a small ground
set, you can choose the large sets carefully --- make them consecutive intervals ---
so that every intersection is itself a consecutive interval, which is automatically
an AP. With a larger ground set, you run out of luck: there are too many possible
triples, and some intersection will inevitably land on a non-AP.

The numbers $5$ and $6$ are the largest $N$ for which the ground set is small enough
to allow this careful arrangement. At $N=7$, the Bron--Kerbosch computation confirms
that the maximum family size drops back to $\binom{7}{2}+1 = 22$, with no exceptional
gain.

There is a faint but pleasing echo here of other classification phenomena in
mathematics. The exceptional Lie algebras $G_2$, $F_4$, $E_6$, $E_7$, $E_8$ exist
only in specific small dimensions. The sporadic finite simple groups appear as
one-off exceptions to the infinite families. Our exceptional values $N=5,6$ are
humbler creatures, but they arise from the same basic mechanism: a structural
argument that works generically fails in a handful of small cases where the
constraints are loose enough to permit something extra.

\section{Where the Formal Proof Stands}

The Lean~4 formalization stands at an interesting inflection point. The lower
bound is fully proved: twelve theorems (\texttt{ap\_intersecting\_star\_N}
and \texttt{ap\_intersecting\_exceptional\_N} for each $N \leq 10$),
verified by \texttt{native\_decide} with zero human case analysis. The upper
bound is proved on paper, but its formalization depends on five axioms that
remain to be discharged:

\begin{center}
\small
\begin{tabular}{ccl}
\toprule
\# & Axiom & What It Says \\
\midrule
1 & \texttt{upper\_bound\_small\_N}
    & $N \leq 10$: no AP-intersecting family exceeds
      $\operatorname{maxSize}(N)$ \\
2 & \texttt{exists\_maximal\_extension}
    & Any family extends to an inclusion-maximal one \\
3 & \texttt{star\_reduction}
    & A maximal family transforms into a star without shrinking \\
4 & \texttt{star\_size\_bound}
    & $N \geq 7$: star AP-intersecting families have no sets
      of size $\geq 4$ \\
5 & \texttt{star\_max\_size}
    & A restricted star has at most $\CN + 1$ members \\
\bottomrule
\end{tabular}
\end{center}

Axioms~2 and~5 are straightforward: the first follows from finiteness of the power
set, the second is a direct counting argument. Axiom~1 is a finite computation ---
the Bron--Kerbosch output serves as a certificate, and formalizing that certificate
in Lean is a matter of engineering. Axioms~3 and~4 are the mathematical core: the
star reduction (shifting/compression) and the size bound (case analysis on triples).
Closing them would yield a complete, machine-checked proof of the theorem.

The current Lean source sits at 304 lines with zero sorries on the lower-bound side.
\end{document}
